\section{Kernel effects and the named-interface filtration}\label{sec:examples}

\subsection{Kernel amplification and kernel obstruction}

The next two examples show why homomorphic coherence alone does not describe incomplete rows.

\begin{example}[One supplied cell forces the whole row]\label{ex:amplification}
Let $A=\{0,1,2\}$ with one binary operation $f$ given by
\[
\begin{array}{c|ccc}
f&0&1&2\\\hline
0&0&0&0\\
1&0&0&0\\
2&0&1&0
\end{array}
\]
and take one operator $b$.  Supply only
\[
\alpha(b,2)=0.
\]
The generated input subalgebra is $S=\{0,2\}$.  Compatibility gives
\[
\alpha(b,0)=f(\alpha(b,2),\alpha(b,2))=0.
\]
The induced map $h:S\to A$ is therefore constant zero, so its kernel identifies $0$ and $2$.  But
\[
\Cg_A(0,2)=\nabla_A,
\]
because the context $f(-,1)$ sends $0$ and $2$ to $0$ and $1$.  Hence Theorem~\ref{thm:kernelsat} gives $D=A$: the single supplied cell forces
\[
\alpha(b,0)=\alpha(b,1)=\alpha(b,2)=0.
\]
The carrier nevertheless remains protected, since the constant-zero map is an endomorphism of $A$.
\end{example}

This is pure determinacy without confusion: kernel propagation enlarges the forced domain while preserving all carrier distinctions.

\begin{example}[A coherent partial row can still force confusion]\label{ex:cepobstruction}
Let $A=\{0,1,2,3\}$ with binary operation
\[
\begin{array}{c|cccc}
f&0&1&2&3\\\hline
0&0&2&0&0\\
1&1&2&1&3\\
2&0&2&2&0\\
3&3&0&3&0
\end{array}
\]
and supply
\[
\alpha(b,0)=0,
\qquad
\alpha(b,2)=2,
\qquad
\alpha(b,3)=0.
\]
Write $r_i=\alpha(b,i)$. The generated domain is $S=\{0,2,3\}$,
and $h(0)=h(3)=0$, $h(2)=2$ is a homomorphism $S\to A$.
Its kernel identifies $0$ and $3$. The context $f(-,1)$ adds $2\equiv0$
to $\Cg_A(0,3)$, while $f(1,-)$ adds $1\equiv3$. Thus
\[
\Cg_A(0,3)=\nabla_A,
\qquad
\Cg_A(\ker h)\cap S^2>\ker h.
\]
The raw row is not protected. Notice that this universal \emph{row}
congruence does not imply universal collapse of the distinguished carrier.
The carrier loses exactly the distinction between $0$ and $2$.
Indeed, compatibility and the supplied cells give
\[
r_2=f(r_0,r_1)=f(r_3,r_1)=r_0,
\]
so $E\vdash 2=0$. The partition
\[
\theta=\{0,2\}\mid\{1\}\mid\{3\}
\]
is a congruence, and the constant row with value $[0,2]$ is an endomorphism
of $A/\theta$ satisfying the supplied data. This quotient model separates
its three classes and proves $\theta_E^A=\theta$.

The missing cell is forced as well. Since $f(1,0)=1$ and $f(1,3)=3$,
\[
r_1=f(r_1,r_0)=f(r_1,r_3)=r_3=0.
\]
Therefore all four named row values coincide with the carrier class
$[0,2]$. No external named row value remains in this example.
\end{example}

Examples~\ref{ex:amplification} and \ref{ex:cepobstruction} separate two effects of the same kernel.  It can spread through $A$ and determine new row values without confusion, or it can feed additional equality back into its original domain and thereby force a carrier quotient.  The latter is the local congruence-extension obstruction.


\begin{example}[External values can coexist with a completion]\label{ex:externalcompletion}
Let $A=\{0,1\}$ with $x\vee y=\max(x,y)$ and one entirely unspecified row.
The named carrier and the row copy are disjoint in the initial algebra by
Theorem~\ref{thm:globalrows}. Both row values are external. Nevertheless,
the identity endomorphism of $A$ is a carrier-valued completion.
A choice of completion does not turn its chosen values into consequences
of the original equations.
\end{example}
